\begin{table}[htbp]
\centering
\small
\caption{Headline quantities of twelve of the sixteen contributions, read from the stored documents of the case named in the second column, each as a count or a point value without an interval; the Monte Carlo errors are printed beside each value in the text of the method, and the four methods of Sections~\ref{sec:bs13} to~\ref{sec:bs16} report theirs in the tables on the complement to the attribution test, the anchored probability of default, the sequential monitors and the conformal recalibration.}
\label{tab:proposed-summary}
\fitwidth{%
\begin{tabular}{lllr}
\toprule
Method & Case & Quantity measured & Value \\
\midrule
Size and power of the tests & C22 & tests of the null family within the size bound at 5\% & 20 of 21 \\
Size and power of the tests & C22 & tests of the null family within the size bound at 5\% and at 1\% & 19 of 21 \\
Reconciliation with ORE & C12 & single swap, dates within two combined standard errors & 33 of 33 \\
Reconciliation with ORE & C12 & CSA with 10 days, dates within two combined standard errors & 17 of 33 \\
Informative withdrawals & C04 & relative error of the B probability of default, withdrawals removed, informativeness 9 (\%) & -34 \\
Thin grades & C04 & coverage of the Wald interval for the first grade at 50 obligors (\%) & 0.33 \\
UF as its own currency & C15 & indexation charge over the change in net interest income, short-UF variant (\%) & 55.9 \\
Dependence band & C18 & worst-case capital over the simple sum & 1.25 \\
Reverse liquidity stress & C16 & run-off multiplier at the 30-day target, adverse base & 2.118175 \\
Reverse liquidity stress & C16 & fast-run base & refused \\
Operational-risk instability & C17 & interquartile range over the truth, splice rung, range over the cells & 1.35 to 3.30 \\
\bottomrule
\end{tabular}}
\end{table}
